% Fragmento integrable. Preambulo anfitrion: amsmath, amssymb, longtable, hyperref.
\section{Compatibilité des lecteurs et restitution du paramètre}
\label{hmtfg:fragmento:compatibilidad-de-lectores-y-restitucion-del-parametro}

\subsection{Compatibilité des lecteurs, des minima et restitution de l’histoire}
\label{hmtfg:escalado:13}

Les opérations de prolongement et les lecteurs ordonnés conservent l’information requise pour la sélection des racines. Les identités suivantes fixent leurs domaines et leur composition.

\subsubsection{Règle forward restituée}
\label{hmtfg:escalado:13.1}

L’actualisation déterministe part de l’état de frontière R36 avec le caractère \(\chi\) et le reste interne. L’actualisation est
\begin{equation}
\mathcal U_K^\chi=
\operatorname{Upd}_K^\chi\circ
\operatorname{Tra}_K^\chi\circ
\operatorname{Sel}_K^\chi,
\qquad
N_{K+1}=729N_K+d_K^\chi.
\label{hmtfg:escalado:eq:26}
\end{equation}
Le chiffre et l’actualisation dépendent de l’état précédent. Deux sections forward ayant le même état complet R36 coïncident par induction. La sous-relation \(\operatorname{Ext}^{\chi,\rightarrow}\) est le graphe de cette actualisation ; la relation totale \(\operatorname{Ext}\) admet une multiplicité de continuations selon les données et les caractères.

Les prolongements compatibles transportent les lecteurs régionaux et la signature conjointe sur les mêmes préfixes.

Ces règles transportent l’orientation, le reste, le quotient, la retenue, la frontière et la mémoire. Leur application à la cascade doit également conserver la sélection ordonnée du paramètre du lecteur critique.

\subsubsection{Indépendance des constantes antécédentes à l’égard du paramètre libre}
\label{hmtfg:escalado:13.2}

Soit \(B\) l’état HMT qui fournit la roue et ses constantes antécédentes. Notons \(\mathcal C(B)\) l’ensemble de ces lectures et \(u_B\) le profil critique. Dans le domaine du lecteur critique, la construction ultérieure est
\begin{equation}
(B,\lambda)\longmapsto f_{B,\lambda}(x)=1-\lambda u_B(x),
\qquad 0<\lambda\le 2/u_B(1).
\label{hmtfg:escalado:eq:27}
\end{equation}
Pour un même \(B\), la lecture des constantes antérieures se factorise par la projection sur la première composante :
\begin{equation}
\widetilde{\mathcal C}(B,\lambda)
=\mathcal C\circ\operatorname{pr}_B(B,\lambda)
=\mathcal C(B).
\label{hmtfg:escalado:eq:28}
\end{equation}
Par conséquent, deux valeurs distinctes de \(\lambda\) conservent les mêmes lectures antécédentes. L’égalité est une propriété de cette construction causale : le profil est fixé avant de faire varier \(\lambda\), et le facteur \(\pi_{\mathrm{HMT}}\) se simplifie dans sa normalisation quadratique.

\textbf{Lemme.} Supposons qu’une condition \(A\) soit évaluée uniquement sur ces lectures antécédentes. Pour tout \(B\) fixé,
\begin{equation}
\{\lambda\in\Lambda:A(\widetilde{\mathcal C}(B,\lambda))\}
=
\begin{cases}
\Lambda,&A(\mathcal C(B))\text{ est satisfaite},\\
\varnothing,&A(\mathcal C(B))\text{ échoue}.
\end{cases}
\label{hmtfg:escalado:eq:29}
\end{equation}
\textbf{Démonstration.} La substitution de \eqref{hmtfg:escalado:eq:28} rend le prédicat indépendant de \(\lambda\). Par conséquent, tous les paramètres du domaine le satisfont, ou aucun. Cela inclut une comparaison des constantes antécédentes avec des intervalles métrologiques.

L’unicité de la section \(B\) et la restriction ultérieure \(f_{B,\lambda}^{2^n}(0)=0\) produisent une fibre de solutions dans \(\lambda\). La règle de première racine agit sur cette fibre par le retour et son itinéraire. La sélection et son prolongement sont démontrés dans les résultats suivants.

Ce résultat maintient l’ordre HMT : la métrologie confronte des sorties préalablement construites et la sélection de la racine se démontre dans le lecteur qui la produit.

\subsubsection{Lemme de transport ordonné du sélecteur}
\label{hmtfg:escalado:13.3}

Soit \(P_n\) un ensemble ordonné de candidats généalogiques de minimum \(p_n^*\). Soit \(Z_n\) l’ensemble initial des racines candidates du corpus, muni de l’ordre réel, et soit \(L_n:P_n\to Z_n\) un lecteur monotone. Son image est dite \textbf{coinitiale} lorsque
\begin{equation}
\forall z\in Z_n\quad\exists p\in P_n:
L_n(p)\le z.
\label{hmtfg:escalado:eq:30}
\end{equation}
Alors
\begin{equation}
\boxed{L_n(p_n^*)=\min Z_n
\quad\Longleftrightarrow\quad
L_n(P_n)\text{ est coinitial dans }Z_n.}
\label{hmtfg:escalado:eq:31}
\end{equation}
\textbf{Démonstration.} Si l’image du minimum est le minimum global, on prend \(p=p_n^*\) dans \eqref{hmtfg:escalado:eq:30}. Réciproquement, étant donné \(z\in Z_n\), la coinitialité fournit \(p\) avec \(L_n(p)\le z\). La monotonie donne \(L_n(p_n^*)\le L_n(p)\le z\). Puisque \(L_n(p_n^*)\in Z_n\), c’est le minimum global. La surjectivité de \(L_n\) constitue une condition suffisante, plus forte que \eqref{hmtfg:escalado:eq:30}.

L’isomorphisme ordonné de la carte du continu satisfait la surjectivité du lecteur. L’application concrète au retour et à son ensemble complet de zéros est démontrée dans la section de transport ordonné ; la classification et l’isolement déterminent ensuite sa continuation dynamique.

\subsubsection{Sélection minimale par survie à profondeur arbitraire}
\label{hmtfg:escalado:13.6}

La composition avec la survie produit également un sélecteur global précis. On conserve la famille uniforme \(\eta=0\), construite dans la section~\ref{hmtfg:escalado:1}, et l’on utilise les théorèmes de réalisation de l’itinéraire infini et de conservation de ses signes. Soient
\[
I=\left[\frac1{2u_0(1)},\frac2{u_0(1)}\right],\qquad
c_j(\lambda)=f_\lambda^j(0),\qquad
\tau_j=(-1)^{v_2(j)}.
\]
L’ensemble
\[
\Lambda_\tau=\{\lambda\in I:\operatorname{sign}c_j(\lambda)=\tau_j
\text{ pour tout }j\ge1\}
\]
est compact et non vide. Le contrôle des dérivées de cette source fournit
\begin{equation}
B_j=\frac{16^j(16^j-1)}{15},\qquad
|c_j(\lambda)|>\frac2{B_j}\quad(\lambda\in\Lambda_\tau).
\label{hmtfg:escalado:eq:32}
\end{equation}
Définissons
\begin{equation}
\varepsilon_j=\frac1{B_j},\qquad
A_N=\{\lambda\in I:\tau_jc_j(\lambda)\ge\varepsilon_j,\
1\le j\le N\}.
\label{hmtfg:escalado:eq:33}
\end{equation}
Chaque \(A_N\) est compact par continuité des itérées et non vide parce qu’il contient \(\Lambda_\tau\). De plus,
\begin{equation}
A_{N+1}\subseteq A_N,\qquad
\bigcap_{N\ge1}A_N=\Lambda_\tau.
\label{hmtfg:escalado:eq:34}
\end{equation}
L’inclusion de droite à gauche découle de \eqref{hmtfg:escalado:eq:32}. L’inclusion inverse découle des marges strictement positives de \eqref{hmtfg:escalado:eq:33}, qui fixent chaque signe.

\textbf{Théorème du minimum survivant.} Les minima \(a_N=\min A_N\) satisfont
\begin{equation}
\boxed{a_N\uparrow a_*=\min\Lambda_\tau.}
\label{hmtfg:escalado:eq:35}
\end{equation}
\textbf{Démonstration.} L’inclusion \eqref{hmtfg:escalado:eq:34} rend croissante la suite des minima, bornée dans \(I\) ; elle admet donc une limite \(a\). Pour chaque \(r\), tous les termes \(a_N\), \(N\ge r\), appartiennent au fermé \(A_r\) ; ainsi \(a\in A_r\). Par \eqref{hmtfg:escalado:eq:34}, \(a\in\Lambda_\tau\). Pour tout \(\lambda\in\Lambda_\tau\), on a \(a_N\le\lambda\) à chaque niveau et, en passant à la limite, \(a\le\lambda\). Cela prouve \eqref{hmtfg:escalado:eq:35}.

Dans le lecteur cylindrique HMT, l’objet \(a_*\) possède des préfixes compatibles à toute profondeur, avec conservation de l’orientation, du reste et de la mémoire. Son choix utilise l’ordre de la lecture paramétrique et la survie complète. La compacité démontre l’existence et la convergence de ces minima.

Les paramètres \(a_N\) résolvent des contraintes de signes avec marge. Leurs frontières peuvent satisfaire \(\tau_jc_j=\varepsilon_j\), tandis que les racines superstables satisfont \(c_{2^n}=0\). Ainsi, \eqref{hmtfg:escalado:eq:35} établit un sélecteur global de réalisation infinie et préserve, comme objet distinct, le sélecteur de premières racines du corpus.

\subsubsection{Restitution du paramètre à partir de l’histoire critique complète}
\label{hmtfg:escalado:13.7}

Le deuxième terme de l’histoire critique fournit un lecteur inverse exact :
\begin{equation}
c_1(\lambda)=1,\qquad
c_2(\lambda)=1-\lambda u_0(1),\qquad
\boxed{\lambda=\frac{1-c_2(\lambda)}{u_0(1)}.}
\label{hmtfg:escalado:eq:36}
\end{equation}
La positivité \(u_0(1)>0\), démontrée pour le profil, rend la division valide. Par conséquent, deux histoires critiques complètes identiques ont le même paramètre. Le mot de signes \(\tau\) constitue une lecture réduite de cette histoire ; il conserve sa combinatoire et omet la grandeur de \(c_2\).

Les formules \eqref{hmtfg:escalado:eq:35}–\eqref{hmtfg:escalado:eq:36} unissent la sélection par survie et la restitution du paramètre à partir de l’histoire critique. L’égalité de sa limite avec le croisement stable est démontrée par première sortie, et l’identification des centres par le transport analytique de leurs classes hybrides.
